\section{The retained instrument realization theory and its boundary}
\label{sec:r74-instruments}

The group results concern terminal responses. For an outcome at every call,
one must instead preserve each command--outcome history and every subsequent
decision made from that visible history. This stronger interface is already
covered, under a sufficient uniform posterior condition, by the earlier
instrument theory of this paper. We recall its exact scope before proving
the new conditional-memory classification.

For fixed known completely positive branches $\mathcal I_{a,y}$ whose
command sums are trace preserving, the hypothesis is
\[
 \mathcal I_{a,y}(X)\succeq
   \frac\eta h\tr\bigl(\mathcal I_{a,y}(X)\bigr)I
 \quad(X\succeq0),\qquad \eta>0.
\]
It is a condition on every recorded branch, including the zero-branch
convention, and follows when depolarizing noise of weight $\eta$ is applied
after each branch. In every finite dimension $h\ge2$, and for $k$ prescribed
preparations, Theorem~\ref{thm:r68-instrument} gives at most
\[
 k+\left\lceil(4h/\eta)^{h-1}\right\rceil
\]
hidden labels, with at most $h^2$ successors per branch row. All commands use
the same physical dictionary. The complete proof in
Section~\ref{sec:interior-instruments} gives the operator identity for every
history; it preserves all adaptive stopped laws, including nontermination,
and nonnegative costs. The hidden label is not exposed to the observer.
Corollary~\ref{cor:r68-risk} consequently preserves the risk of every common
report-based observer, including constrained boundary budgets; it does not
identify hidden simulator size with optimal observer memory.

The same retained theory includes a terminating algebraic construction
(Proposition~\ref{prop:r68-algebraic}), interiority after a finite block with
the clock and intermediate reports counted (Theorem~\ref{thm:r68-block}),
and matching powers of $\eta^{-1}$ for universal state-positive
representation (Theorem~\ref{thm:r68-noise-exponent}). The latter lower bound
is for generators which are actual density matrices. It is not imported
as a lower bound against arbitrary nonnegative hidden-state representations.
The family in Section~\ref{sec:r74-mixing-obstruction} will have a lower
bound against that unrestricted class.

Invariant positive cones and their stochastic realizations are classical
\cite{Heller65,MonrasWinter14}. Okay's quantum simulation framework explicitly
includes invariant-polytope updates and adaptive composition
\cite{Okay26}. Pure-state nets and approximation of the density-matrix body
have an established convex-geometric theory \cite{AubrunSzarek17}; extension
complexity is a different resource \cite{Fawzi21}. These ingredients, and
the paper's already proved interior-instrument theorem, are not asserted
as new discoveries of this revision. Their role is to fix the sufficient
condition and the resource convention against which the new sharp examples
are compared. The next common-contraction construction gives an explicit
certificate in supplied linear-response coordinates; it is also an
application of positive-realization geometry, not a replacement for the
retained instrument theorem.
